{\small
\begin{longtable}{@{}>{\raggedright\arraybackslash}p{0.27\textwidth}>{\raggedright\arraybackslash}p{0.68\textwidth}@{}}
\caption{Thirty-four cases, grouped by failure mode, with the cause in the words of the primary record: a retraction notice, a finding of the Office of Research Integrity, or a published reanalysis.}
\label{tab:fc-cases} \\
\toprule
\textbf{Case} & \textbf{What its primary record says} \\
\midrule
\endfirsthead
\toprule
\textbf{Case} & \textbf{What its primary record says} \\
\midrule
\endhead
\bottomrule
\endfoot
\multicolumn{2}{@{}l}{\textbf{Printed value disagrees with the authors' own records}} \\ \nopagebreak
Abegaz et al. (2014), PLOS ONE & ``the estimated number of deaths and injuries reported cannot be reproduced when using the unmatched and matched police and hospital data''~\citep{plosone2026abegaz} \\
Sun et al. (2024), PLOS ONE & ``the p50 median value is lower than the reported minimum value''~\citep{plosone2026sun}. Also states that data, images or records could not be produced. \\
Wansink et al. (2012), Archives of Pediatrics and Adolescent Medicine & ``We had erroneously reported the age group as children ranging from 8 to 11 years old; however, the children were 3 to 5 years old.''~\citep{wansink2017retraction} \\
\midrule
\multicolumn{2}{@{}l}{\textbf{Report departs from the registered plan}} \\ \nopagebreak
67 trials in five medical journals, 2015 to 2016 (a cohort, not one publication) & ``365 novel outcomes were reported without declaration in total, a mean of 5.4 per trial''~\citep{goldacre2019compare} \\
\midrule
\multicolumn{2}{@{}l}{\textbf{References that do not exist}} \\ \nopagebreak
Khalilzadeh et al. (2026), PLOS ONE & ``Citation of unretrievable references 1, 2, 4, 5, 8, 10, 12--18, 20--22, 33, 36, 51, 61, 75, and 79.''~\citep{plosone2026khalilzadeh} \\
Mohammad et al. (2026), PLOS ONE & ``Citation of unretrievable References 2, 3, 6, 7, 8, 18, 20, and 26.''~\citep{plosone2026mohammad} \\
\midrule
\multicolumn{2}{@{}l}{\textbf{Data could not be produced for audit}} \\ \nopagebreak
Mehra et al. (2020), Lancet & ``Surgisphere would not transfer the full dataset, client contracts, and the full ISO audit report to their servers for analysis''~\citep{mehra2020retractionlancet} \\
Mehra et al. (2020), New England Journal of Medicine & ``the raw data could not be made available to a third-party auditor, we are unable to validate the primary data sources''~\citep{mehra2020retractionnejm} \\
\midrule
\multicolumn{2}{@{}l}{\textbf{Error inside the analysis}} \\ \nopagebreak
Morisky et al. (2008), Journal of Clinical Hypertension & an independent statistical review ``concluded that the results were misleading due to issues regarding the sensitivity and specificity of the medical adherence scale used''~\citep{wiley2023morisky} \\
Gebrechorkos et al. (2025), Nature & ``area-weighted grid cells were not used when calculating study-wide averages, an error that affected numerical values in the main text''~\citep{gebrechorkos2026retraction} \\
Islam et al. (2019), PLOS ONE & ``an error in the p-value calculation caused by a coding error''~\citep{plosone2026islam} \\
Zhang et al. (2025), PLOS ONE & a formula ``contains an incorrect minus sign, which should be a plus sign''~\citep{plosone2026zhang} \\
Okray et al. (2023), Nature & the authors found ``errors in the pipeline originally used for data analysis''~\citep{okray2026retraction} \\
Hewitson et al. (2021), PLOS ONE & the correlation analysis ``included typically developing (TD) participants in error''~\citep{plosone2024hewitson} \\
Rahman et al. (2021), PLOS ONE & ``the PSS mean reported in [1] is consistent with incorrect coding of these items''~\citep{plosone2026rahman}. Also states that data, images or records could not be produced. \\
Huang et al. (2021), Lancet & an expression of concern was published ``because of the substantive data correction request from the authors following comments from a reader''~\citep{lancet2023huang} \\
Reinhart and Rogoff (2010), American Economic Review & ``coding errors, selective exclusion of available data, and unconventional weighting of summary statistics lead to serious errors''~\citep{herndon2013debt} \\
\midrule
\multicolumn{2}{@{}l}{\textbf{Data fabricated or falsified}} \\ \nopagebreak
Potti et al. (2006), Nature Medicine & ``Respondent altered data sets to improve the accuracy of predictors for response to treatments''~\citep{ori2015potti}; ``the results incorporate several simple errors that may be putting patients at risk''~\citep{baggerly2009deriving}; ``a corruption of several validation data sets precludes conclusions regarding these signatures''~\citep{potti2011retraction} \\
Potts Kant et al. (2013), Pediatric Research & falsified data by ``changing numerical values in the datasets generated by the apparatus''~\citep{ori2019pottskant} \\
Savine and Braver (2012), Cognitive, Affective and Behavioral Neuroscience & ``Respondent exaggerated (1) the effect of incentive context on response times and error rates in Table 1''~\citep{ori2013savine} \\
Skau et al. (2016), Cell & falsification by ``selectively reporting by inappropriate inclusion/omission or alteration of data points in ten (10) figures''~\citep{ori2018skau} \\
Hauser et al. (2002), Cognition & ``Respondent published fabricated data in Figure 2 of the paper''~\citep{ori2012hauser} \\
Linn et al. (2003), Cellular and Molecular Biology & ``he falsified the values in Tables 2-5 by altering the values that he had obtained from another author's manuscript''~\citep{ori2010linn} \\
Thiruchelvam et al. (2005), Journal of Biological Chemistry; Rodriguez et al. (2005), Environmental Health Perspectives & misconduct by ``falsifying and fabricating cell count data that she claimed to have obtained through stereological methods''~\citep{ori2012thiruchelvam} \\
Taylor and Gercel-Taylor (2008), Gynecologic Oncology & two investigations ``found evidence of falsification, fabrication, and plagiarism of data''~\citep{taylor_notice_2023} \\
\midrule
\multicolumn{2}{@{}l}{\textbf{Image manipulation or duplication}} \\ \nopagebreak
Jiang et al. (2002), Nature & ``The original images for Figures 6a, 6e and 6m could not be retrieved by the authors''~\citep{jiang2024retraction} \\
Lesné et al. (2006), Nature & the figures ``show signs of excessive manipulation, including splicing, duplication and the use of an eraser tool. The data cannot be verified from the records.''~\citep{lesne2024retraction} \\
Zou et al. (2001), Nature & falsified figures by ``manipulating the images to alter the number and location of positively stained cells''~\citep{ori2014zou} \\
Liu et al. (2014), PLOS ONE & ``the data underlying the panels of concern are no longer available''~\citep{plosone2026liu} \\
\midrule
\multicolumn{2}{@{}l}{\textbf{Study conducted differently from its description}} \\ \nopagebreak
Wakefield et al. (1998), Lancet & the paper's claims ``that children were `consecutively referred' and that investigations were `approved' by the local ethics committee have been proven to be false''~\citep{lancet2010wakefield} \\
Estruch et al. (2013), New England Journal of Medicine & ``Because of irregularities in the randomization procedures, we wish to retract the following article''~\citep{estruch2018retraction} \\
Gautret et al. (2020), International Journal of Antimicrobial Agents & ``The journal has not been able to establish whether all patients could have entered into the study in time''~\citep{gautret2025retraction}. Also states that data, images or records could not be produced. \\
LaCour and Green (2014), Science & ``Survey incentives were misrepresented.''; ``LaCour has not produced the original survey data from which someone else could independently confirm the validity of the reported findings''~\citep{mcnutt2015retraction} \\
\midrule
\multicolumn{2}{@{}l}{\textbf{None of the above}} \\ \nopagebreak
Fukuhara et al. (2005), Science & ``we have found that not all preparations of visfatin bind to and activate the insulin receptor''~\citep{fukuhara2007retraction} \\
\end{longtable}
}
